% ============================================================
% 第七章 · 最终显示效果
% 对应大纲：# 七、最终显示效果
%
% 本章是成品展示章，以图为主、文字为辅，12 张图，按玩家的实际游玩
% 顺序排列：进入与开场 → 剧情与探索 → 玩法与终局 → 系统与视听。
% 排列依据是游玩顺序而不是功能模块：这 12 张连着看一遍等于把游戏走一遍。
%
% 取舍口径：图控制在 12 张以内，被砍掉的界面由正文文字补上（例如四个
% 小游戏只放一张图，另外三个在导语里点名）。删图不等于删信息。
%
% 图片放 figures/07-final-showcase/，命名见 figures/README.md
% （全小写英文 + 连字符）。
%
% 截图未到位时用 \figph 占位（定义在 tex/commands.tex），占位框不依赖
% 任何图片文件，编译能过。图片就位后把 \figph 改成 \fig、把第二个参数
% 换成图片路径，后三个参数原样保留。占位框里那行小字是给截图的人看的
% 操作说明，换成 \fig 之后自然消失；图注是永久的，已按最终效果写好。
%
% ---------- 12 张截图 ----------
%   7.1  入口页首屏 / 祠堂首屏
%   7.2  人物对白与立绘 / 村口 / 陈家老宅 / 线索详情卡
%   7.3  逃离数据机房 / 证据处置界面 / 五个结局拼图
%   7.4  成就页 / 移动端主游戏页 / 五种转场手法拼图
% 两张需要自己拼图：五个结局的画面拼成一张，五种转场的过程帧拼成一张。
%
% 被砍掉、想加回来时按顺序优先补这几张：制作组页（团队展示）、主菜单、
% 背包（物品与证物两个页签）、序幕视频「白灯初现」的截帧、项目介绍页、
% 另外三个小游戏、注册与登录页、18 个场景的缩略图总览。
%
% ---------- 名称以代码与 V3 文档为准，不用描述性说法 ----------
% 五个结局的正式名称是《共犯的终点》《封井之人》《无名者》《白灯之后》
% 《不再借灯》，出自 docs/docs-v3/V3剧情Node推进.md 与代码里的
% ending-* 节点。报告此前只用「交出证据」「销毁全部」这类触发条件指代
% 它们，现已按正式名统一到第一、二、三、七章。
% 四个小游戏在界面上的标题是「复原手绘地图」「拆解控制室网络」
% 「还原矿井路线」「逃离数据机房」，见 v3-minigame-gateway.js 的注册表。
% 第二章的「拆鬼」「路线解谜」是设计阶段的叫法，不是玩家看到的标题，
% 现已统一。各人个人总结里的措辞是本人原文，未改动。
% ============================================================

\chapter{最终显示效果}

本章按玩家的实际游玩顺序排列成品截图：从入口页进入，经序章、探索与四个小游戏，走到五个结局，最后是系统页面与多端适配。这 12 张连着看一遍，等于把《\ProjectName》走了一遍。

前面几章分别交代了设计规范、剧情结构与技术实现，本章只呈现跑起来之后的样子。截图均在正式版本上取得，不含原型稿与线框图。

\section{进入与开场}

玩家从入口页进入，经登录或游客通道到主菜单，再走进祠堂，游戏正式开始。这一节的两张图分别是全篇第一印象与第一幕。

入口页是全篇唯一使用毛笔手写字体的位置。冷调雨幕压满整屏，白灯是画面上唯一的暖光源，第三章定下的冷暖对比在这里第一次出现。

% 【必】
\figph[0.85]{入口页首屏：整屏雨幕 + 白灯 + 毛笔标题 + 「进入」按钮。浏览器全屏截图，不要带上书签栏和地址栏}{入口页首屏}{entry}

祠堂是序章的全部场景。主角从地上醒来，身上只有三件物品，旁边有一个自称小周的年轻人。这一屏同时是场景层、阅读层与对话框的第一次合流。

% 【必】
\figph[0.85]{祠堂首屏：主角醒来。截白灯亮起后的那一屏（祠堂有三张背景图：只有灯泡、灯泡与灯笼都亮、只有白灯笼），画面里要有小周和几处可点的热点}{祠堂首屏}{prologue-hall}

\section{剧情与探索}

剧情文本分旁白、独白与角色对白三层，渲染上靠名牌区分：旁白不出名牌，独白出名牌，玩家始终知道一句话是谁说的，或者知道它没有来源。立绘只有左右两个站位，站位变化即表示说话人更换。

% 【必】
\figph[0.85]{人物对白：带立绘的一屏，角色站在画面左侧或右侧，名牌与正文分行显示}{人物对白与立绘}{dialogue-npc}

探索以整张场景图作交互载体，全篇 18 个场景，覆盖村庄中枢、三条身份线、三处玩法场景和五个结局画面。场景里可调查的对象一律做成热点，调查后留下痕迹或标记为已调查。下面两张是其中的两个地址：村口是唯一的中枢，玩家在这里反复出入，把三条线的进度带回同一个地方；陈家老宅是陈晋年线的第一站，也是全篇第一处门后主光源。

% 【必】
\figph[0.85]{村口：中枢场景，画面上有多个可点热点。截玩家刚到这里、还没有调查过的状态}{村口·村庄中枢}{scene-village}

% 【必】
\figph[0.85]{陈家老宅：用刻名钥匙开门后的内景，带门后主光源}{陈家老宅}{scene-chen-house}

调查得到的物证进入背包，线索独立成页。玩家拿到的是物证，理解的是线索，证据链由线索之间互相印证构成。

% 【必】
\figph[0.85]{线索详情卡：点开一条关键线索后的卡片，朱砂红只出现在关键真相上}{线索详情卡}{clue-card}

\section{玩法与终局}

四个小游戏分别嵌在剧情节点上，玩完把结果回写给剧情模块，由剧情模块决定后续走向。复原手绘地图在村口，玩家从三名村民处取得碎片拼成村庄图；拆解控制室网络在广播室，逆向追踪童谣、白灯与湿脚印三类人工装置；还原矿井路线在旧矿井，从排水洞绕过封堵进入竖井；逃离数据机房是终局的俯视追逐，躲避小周的追捕、收集三份证据并撤离。下面这张是第四个，它在其中操作要素最全。

% 【必】
\figph[0.85]{逃离数据机房：俯视视角的追逐场面，画面里要有生命值、三份证据的收集进度和追捕者小周}{逃离数据机房与生命值}{minigame-chase}

终局最后一步不交给对话，而是给玩家一个证据处置界面。三个选项分别处理手里的材料：销毁全部证据、以白灯客的名义选择性公开、完整公开。加上对峙时的交出证据与追逐失败，全篇共五个稳定结局，每个结局配一段过场视频。

% 【必】
\figph[0.85]{证据处置界面：村口场景上的三个处置选项，是整个游戏的最终抉择}{证据处置界面}{evidence-disposition}

% 【必】需要自己拼图：五张结局画面横排拼成一张（或两行，上二下三）
\figph[0.85]{五个结局的画面拼成一张，从左到右依次是《共犯的终点》（交出证据）、《封井之人》（追逐失败）、《无名者》（销毁全部）、《白灯之后》（选择性公开）、《不再借灯》（完整公开）。需要自己拼图}{五个结局}{endings}

\section{系统与视听}

主菜单、成就页与成员页在主线之外，玩家中途可以随时返回。成就页收录 14 条成就，按剧情推进分成调查札记、灯下真相、终局余烬三组；五条结局成就被标记为隐藏条目，未解锁时只显示未明之章配合罗马数字序号，不渲染描述文案。

% 【必】
\figph[0.85]{成就页「雨夜留痕」：顶部解锁进度与进度条，下面是三组共 14 张卡片。截已解锁与未解锁混合的一屏}{成就页}{achievements}

全站按窄屏重排，热点与阅读框在 390 像素宽下依然可点可读。

% 【必】
\figph[0.22]{手机端竖版主游戏页：窄屏下的场景、阅读框与底部操作栏。用真机或开发者工具的手机模式截}{移动端主游戏页}{mobile-game}

转场是全篇唯一使用高强度效果的位置，共十一处，分五种手法：轻闪黑、窒息式闪黑、暴力抖动、无信号花屏、抖动接黑。日常动效压到几乎不可察觉，效果强度集中在这里。

% 【必】需要自己拼图：五种手法各取开始、高潮、结束三帧，横排五行
\figph[0.85]{五种转场各取过程帧拼成一张，每种取开始、高潮、结束三帧。窒息式闪黑与无信号花屏最出效果，至少这两种要拼进去}{五种转场手法}{transition-frames}
